\subsection{An evaluated one-place class-two correction for the same ordinary ghost}

\begin{statement}{A.210}{Proved}
{The order-eight quotient and the exact odd correction coset at 1129}
Let \((A,B)\) be the same completed pair constructed in A.207,
and put \(p=1129\).  Every continuous global class-two lift
of this pair has, at \(p\), the same image as the rational
point \(3\) in an actual quotient with geometric group
\(D_8\) of order eight.  Its ordinary point image at
\(t_p\in828+p\Z_p\) has a different image in this quotient.
Consequently no global class-two lift of this pair matches
its ordinary point image even at this one finite place.

More precisely, align the actual coefficient-two abelian
cocycles of any such global lift with those of \(t_p\).
The resulting local central difference of A.40 is
\[
 \delta_{p,2}\in H^1(\Q_p,\Z_2(2))
       =\Z_2/(p^2-1)\Z_2=\Z/16,
                 \qquad\delta_{p,2}\equiv1\pmod2.
\tag{A.210.1}
\]
Across the twenty-four global class-two section classes
above \((A,B)\), the differences are exactly the eight
odd elements of \(\Z/16\), each occurring three times.
The actual global correction map is
\[
 \Z/24\longrightarrow\Z/16,\qquad
 k\longmapsto\frac{1129^2-1}{24}k
                     =53110k\equiv6k\pmod {16}.
\tag{A.210.2}
\]
Thus this ordinary pair is locally matching at the full
abelian level and first fails locally matching continuous
nilpotent lifting at class two.  Its actual global
class-four lift from A.207 remains in place.
\end{statement}

\begin{proof}
\emph{The actual class-two gauge and its dihedral quotient.}
Use the fixed rational tangential splitting and geometric
generators of A.43.  Work first in the coefficient-two
geometric class-two group, with \(C=[x,y]\) and coordinates
\(x^a y^b C^c\), \(a,b,c\in\Z_2\).  The group law is
\[
 (a,b,c)(a',b',c')=(a+a',b+b',c+c'-a'b).
\tag{A.210.3}
\]
The actual Ihara action is pure in this quotient:
\(g(x)=x^{\chi_2(g)}\), \(g(y)=y^{\chi_2(g)}\).
Indeed its conjugating element \(f_g\) lies in
\(\Gamma_2\), which is central here.  In particular
\(g(C)=C^{\chi_2(g)^2}\).

Define the finite group and the literal quotient by
\[
 D_8=\langle r,s:r^4=s^2=1, srs=r^{-1}\rangle,
 \qquad x\longmapsto s,\quad y\longmapsto r,\quad
                         C\longmapsto r^2.
\tag{A.210.4}
\]
The image of \(x^a y^bC^c\) is \(s^a r^{b+2c}\).
Its kernel is exactly
\(a\in2\Z_2\), \(b+2c\in4\Z_2\).
This kernel is preserved by the arithmetic action:
for odd \(\chi\),
\(\chi b+2\chi^2c=\chi(b+2c)+2c(\chi^2-\chi)\),
whose second summand is divisible by four.
The map therefore defines an actual arithmetic quotient
of \(E_2\).  Its central commutator subgroup is
\(\langle r^2\rangle\simeq\Z/2\), on which the
arithmetic action is trivial.

\emph{Exact central inertia and the global quadratic difference.}
At any physical prime \(q>3\), take the pro-two tame
generator \(T_q\) with \(\kappa_2(q)(T_q)=1\) and
arithmetic Frobenius \(\Phi_q\).  In the fixed tangential
gauge, write the geometric coordinates of a global
class-two section as
\[
 V=s(T_q)_{\rm geom}=x^m y^n C^c,\qquad
 U=s(\Phi_q)_{\rm geom}=x^{A_q}y^{B_q}C^{H_q}.
\tag{A.210.5}
\]
The action is unramified at \(q\).  Its wild inertia
has zero image in this pro-two group, and its pro-two
tame image satisfies the literal relation
\(U\varphi_q(V)U^{-1}=V^q\).
The exact multiplication (A.210.3) gives
\[
 \begin{aligned}
 U\varphi_q(V)U^{-1}
   &=x^{qm}y^{qn}C^{q^2c+q(A_qn-B_qm)},\\
 V^q&=x^{qm}y^{qn}C^{qc-\binom q2mn},\\
 c&=\frac{B_qm-A_qn}{q-1}-\frac{mn}{2}.
 \end{aligned}
\tag{A.210.6}
\]
The last equality is initially in \(\Q_2\), with its
left side already integral.  It is in the same tangential
gauge as the Ihara action; no change to the point-three
splitting has been inserted into this calculation.

Outside the new-prime supports of A.207, the two inertia
coordinates are \(m=n=0\), so (A.210.6) gives \(c=0\).
At a new prime, \((m,n)\) is one of
\((E,0),(0,E),(-E,-E)\), where
\[
 E=24n_j(q^2-1)(q^3-1),\qquad
          \frac{E}{q-1}=24n_j(q+1)(q^3-1)\in2\Z.
\tag{A.210.7}
\]
Here \(E\) is divisible by four and \(E^2/2\) is even.
Since \(A_q,B_q\) are integral, (A.210.6) proves
\(c\in2\Z_2\) in all three cases.
Equation (A.210.4) therefore sends the entire inertia
image to one: \(m\) is even and \(n+2c\) is divisible
by four.  The same calculation for the point \(3\),
whose two valuations are zero at every \(q>3\), gives
its trivial dihedral inertia there as well.

The explicit products of A.207 have
\(A=3A_0^{24}\), \(B=-2B_0^{24}\) in the full power
completion.  Choose Kummer representatives using these
actual completed roots.  Geometric conjugation aligns
the abelian representatives of any global lift with them.
The dihedral abelian cocycles then agree literally with
those of the point-three section, modulo two.
Their difference is consequently an actual central
one-cocycle
\[
                \rho:G_\Q\longrightarrow\langle r^2\rangle
                                      \simeq\Z/2.
\tag{A.210.8}
\]
The centre has trivial arithmetic action, so this is a
quadratic character.  The just-proved inertia calculation
shows that it is unramified outside \(2,3,\infty\).
By the literal quadratic Kummer map it is represented by
one of
\[
                           1,-1,2,-2,3,-3,6,-6.
\tag{A.210.9}
\]
Indeed every odd-prime valuation outside three must be
even and is removed by a rational square; the sign and
the valuation parities at two and three give this list.

All eight characters restrict trivially to \(G_{\Q_p}\)
for \(p=1129\).  For example the three generators have
the explicit nonzero square witnesses
\[
                -1\equiv168^2,\qquad
                 2\equiv406^2,\qquad
                 3\equiv365^2\pmod {1129}.
\tag{A.210.10}
\]
Simple-root Hensel lifting makes them squares in
\(\Q_p\).  Thus every global integral class-two lift
has the point-three dihedral section class at this
one finite place.  This conclusion has used its
integral tame relation, not local point matching at
any other prime.

\emph{The explicit finite cover and all local conjugators.}
The dihedral quotient (A.210.4) has the following
actual rational cover over \(U\):
\[
          t=\left(\frac{w^4-1}{w^4+1}\right)^2.
\tag{A.210.11}
\]
Its geometric deck maps are
\(r(w)=iw\), \(s(w)=1/w\).  They satisfy the exact
relations in (A.210.4) and preserve the displayed
function.  The map has degree eight, hence these eight
distinct maps prove its geometric Galois group.
The zeros above zero have ramification index two;
\(1-t=4w^4/(w^4+1)^2\) gives index four at zero and
infinity above one; and the poles above infinity have
index two.  Removing these fibres gives a finite etale
cover of \(U\).

Its meridians at zero and one are respectively a
reflection and a rotation of order four.  The presentation
\(x^2=y^4=(xy)^2=1\) reduces every word to one of
eight words \(y^i x^j\).  Its eight-element image is
therefore exactly (A.210.4), with \([x,y]=y^2\).
This proves equality of the actual geometric kernels,
not merely equality of the group orders.  The cover
is defined over \(\Q\), in agreement with the
arithmetic stability already proved.

At \(p=1129\), \(i\) belongs to \(\Q_p\).
The point-three fibre of (A.210.11) has the eight
simple roots
\[
                 66,202,325,408,721,804,927,1063
                                      \pmod p.
\tag{A.210.12}
\]
They all lift to \(\Q_p\), and trivialize this fibre
as a full local Galois set.  This is an actual local
arithmetic frame furnished by the point-three section.

For \(t\equiv828\pmod p\), the two square-root residues
of \(t\) are \(244,885\).  Solving the displayed
function for \(w^4\) gives the two inverse values
\[
 \begin{array}{c|c|c|c}
 \sqrt t&w^4=(1+\sqrt t)/(1-\sqrt t)
          &(w^4)^{564}&(w^4)^{282}\\\hline
 244&524&1&-1\\
 885&446&1&-1
 \end{array}\quad\pmod {1129}.
\tag{A.210.13}
\]
Each value is a square and is not a fourth power.
Since the residue field contains the fourth roots of
unity, all four roots of each equation lie in its
quadratic extension, and Frobenius sends every one
of them to its negative.  Thus the actual Frobenius
on all eight sheets is the central element
\(w\mapsto-w=r^2=[x,y]\).
The covers are etale at these fibres, and all roots
and the indicated unit-power classes lift by Hensel.
The same Frobenius conclusion therefore holds on the
whole residue disc, with unramified local action.

The identity action in (A.210.12) and the nonidentity
central action in (A.210.13) cannot be conjugate.
All choices of sheet or geometric conjugator merely
conjugate this regular dihedral action; its two
central elements are unchanged.  This explicitly
accounts for every local finite-sheet ambiguity.
Together with (A.210.8)--(A.210.10), it proves the
one-place mismatch for every global class-two lift.

\emph{The actual continuous central correction.}
Both ordinary coordinates at \(p\) are units.
The pro-two power completion of \(\Z_p^\times\)
is its cyclic two-primary root-of-unity group of
order eight: its principal-unit group is uniquely
two-divisible.  The twenty-fourth-power identities
above therefore give
\(\kappa_2(A_p)=\kappa_2(3)\) and
\(\kappa_2(B_p)=\kappa_2(-2)\) in the full
\(\Z_2(1)\) cohomology, not just modulo two.
Thus the point-three, ordinary-point and global
coefficient-two abelian representatives admit the
actual alignment used in A.40.16.
The vanishing of the local abelian Tate invariants
makes its central comparison class well defined.

The reduction of the central coordinate \(C^h\)
under (A.210.4) is \((r^2)^{h\bmod2}\).
The identity and central Frobenius calculations
therefore give precisely \(\delta_{p,2}\bmod2=1\).
This is unaffected by the preceding alignment:
all its reductions act by geometric conjugation,
already accounted for by the central Frobenius action.
By the actual unramified Frobenius isomorphism of
A.40.13,
\[
 H^1(\Q_p,\Z_2(2))=\Z_2/(p^2-1)\Z_2,
 \quad p^2-1=1274640=16\cdot79665.
\tag{A.210.14}
\]
This proves (A.210.1).

A.40 identifies all global class-two lifts above this
pair as the torsor under \(\Z/24\) generated by
\(\eta_2(g)=(\chi(g)^2-1)/24\).
Its literal Frobenius restriction is (A.210.2).
Its image is the entire even subgroup of \(\Z/16\),
with kernel of order three.  Subtracting it from
any one odd comparison class gives exactly the eight
odd classes, with the asserted multiplicities.
None is zero.  The full abelian section is locally
ordinary by A.207, so its first failure of this
continuous nilpotent local-matching test is exactly
class two, as stated.
\end{proof}

The executed certificate
\path{code/point3_dihedral_1129_certificate.py}, with its
retained output and data, checks the literal Hall signs,
the quotient kernel and its arithmetic stability,
37,485 integral tame-parity instances, the rational cover
and its complete local fibres, all eight supported
quadratic characters, and the exact central correction
map.  The source proves the all-prime inertia and
continuous-cohomology assertions; they are not inferred
from the finite parity sample.
